problem:  
Let \(G\) be a compact, connected Lie group acting smoothly and effectively with cohomogeneity one on a compact, simply connected manifold \(M^n\).  Known classifications (e.g. Grove–Ziller) show that many such \(M\) admit \(G\)‑invariant metrics of nonnegative sectional curvature.  Among these, a particularly interesting and structurally rich class are the **doubly warped cohomogeneity‑one manifolds**—those whose group diagram satisfies \(K_\pm/H \cong \mathbb{S}^{\ell_\pm}\) with both \(\ell_\pm \ge 2\).

**Concrete Open Problem (Positivity Rigidity in a Concrete Family):**  
Within the known families of nonnegatively curved cohomogeneity‑one \(G\)-manifolds with \(K_\pm/H \cong \mathbb{S}^{2}\) and \(G/H\) homogeneous space of rank > 1 (e.g., the Grove–Ziller examples on certain \(S^3\)-bundles over \(S^4\) and their higher‑dimensional analogues), determine whether any of them admit a \(G\)-invariant metric of **strictly positive sectional curvature**.  
In other words, can one ever deform the nonnegatively curved invariant metric to strictly positive curvature **without destroying the orbit structure**?  Or is there a genuine geometric obstruction tied to the group diagram preventing full positivity?

Why is it a “good” problem:  
This refinement isolates a **specific and open frontier** in the study of invariant curvature: the Grove–Ziller construction produces abundant nonnegative curvature by analytic gluing, yet no positive curvature example is known with the same orbit structure.  The problem now asks for an explicit test of whether symmetry and topology enforce a strict gap between nonnegative and positive curvature in this concrete, accessible, and still mysterious family.

It is a "good" problem because:
1. **Grounded in current theory:**  It builds directly on the known classification of cohomogeneity‑one manifolds with nonnegative curvature, zeroing in on explicit families where the existence of positive metrics is unknown but geometrically plausible.
2. **Bridges multiple domains:**  It joins transformation group theory (through the orbit diagram \(G \supset K_\pm \supset H\)), Riemannian analysis (metric deformation and curvature estimates), and topology (structure of \(S^k\)-bundles and characteristic classes).
3. **Simple to state, deep to settle:**  The formulation involves only “Can this concrete nonnegatively curved symmetric manifold be made positively curved?” but a resolution would likely require new deformation or obstruction techniques.
4. **Forward‑looking:**  A positive resolution would yield new examples of strict positive curvature; a negative one would expose a new rigidity mechanism relating curvature, symmetry, and topology—either outcome would advance the field.